%%%%%%%%%%%%%%%%%%%%%%%%%%%%%%%%%%%%%%%%%%%%%%%%%%%%%%%%%%%%%%%%%%%%%%
% Lớp 11 | Giới hạn. Hàm số liên tục | Giới hạn của hàm số
% Dạng 1D3?2-8: Toán thực tế áp dụng giới hạn của hàm số
% Số câu: 29  (sinh tự động bởi tach_cau.py; từ nay thầy sửa trực tiếp file này)
%%%%%%%%%%%%%%%%%%%%%%%%%%%%%%%%%%%%%%%%%%%%%%%%%%%%%%%%%%%%%%%%%%%%%%

%%%%%%%%%%%%%%%% TRẮC NGHIỆM 4 PHƯƠNG ÁN %%%%%%%%%%%%%%%%

% nguồn: Ok_Toan11_TracNghiem_Lan2
\begin{ex}%[1D3H2-8]%[ma:0337a203]
	Một cái hồ đang chứa $400$ m$^3$ nước mặn với nồng độ muối $10$ kg/m$^3$. Người ta ngọt hóa  nước hồ bằng cách bơm nước ngọt vào hồ với tốc độ $3$ m$^3$/phút. Viết biểu thức $C(t)$ biểu thị nồng độ muối trong hồ sau $t$ phút kể từ khi bắt đầu bơm.
	 \choice
		 {$C(t)=\dfrac3{400+3t}$}
		{$C(t)=\dfrac{400+3t}3$} 
		{$C(t)=\dfrac{400+3t}{4\,000}$} 
		{\True  $C(t)=\dfrac{4\,000}{400+3t}$}
	 \loigiai{Khối lượng muối có trong hồ là $400\cdot10=4\,000$ (kg).\\
		Sau $t$ phút lượng nước trong hồ là $400+3t$  (m$^3$).\\
	Nồng độ muối tại thời điểm $t$ phút kể từ khi bơm là $C(t)=\dfrac{4\,000}{400+3t}$ (kg/m$^3$).} 
\end{ex}

% nguồn: Ok_Toan11_TracNghiem_Lan2
\begin{ex}%[1D3H2-8]%[ma:17686da9]
	Trong hồ có chứa $7\,000$ lít nước ngọt. Người ta bơm nước biển có nồng độ muối là $30$ gam/lít vào hồ với tốc độ $15$ lít/phút. Viết biểu thức $C(t)$ biểu thị nồng độ  muối của nước trong hồ sau $t$ phút kể từ khi bắt đầu bơm.
	\choice 
		{$C(t)=\dfrac{700+15t}{30\cdot15t}$} 
		{$C(t)=\dfrac{30\cdot15t}{700+15t}$} 
		 {$C(t)=\dfrac{7\,000+15t}{30\cdot15t}$}
		{\True  $C(t)=\dfrac{30\cdot15t}{7\,000+15t}$}
	\loigiai{Sau  $t$ phút lượng nước trong hồ là $7\,000+15t$ (lít).\\
		Khối lượng muối có trong hồ tại thời điểm $t$ phút kể từ khi bơm là  $30\cdot15t$ (gam).\\
	Nồng độ muối tại thời điểm $t$ phút kể từ khi bơm là $C(t)=\dfrac{30\cdot15t}{7\,000+15t}$ (gam/lít).} 
\end{ex}

% nguồn: Ok_Toan11_TracNghiem_Lan2
\begin{ex}%[1D3V2-8]%[ma:dc7a04a7]
	\immini{Một thấu kính hội tụ có tiêu cự $ f>0$ không đổi. Gọi $d$ và $d'$ lần lượt là khoảng cách từ vật thật và ảnh của nó tới quang tâm $O$ của thấu kính. Ta có công thức $\dfrac1d+\dfrac1{{d'}}=\dfrac1f$  hay $d'=\dfrac{df}{d-f}$. Tìm giới hạn $\lim\limits_{d\to {f^+}}\dfrac{df}{d-f}$.
	 \choice
		 {$-\infty $}
		 {$f^2$}
		{\True $+\infty $} 
		{$f$}}{\begin{tikzpicture}[>=stealth,scale=1,font=\footnotesize]
		\coordinate (O) at  (0,0);
		\coordinate (F') at (1,0); 
		\coordinate (A') at  (3,0);
		\coordinate (a') at (4,0); 
		\coordinate (F)  at (-1,0);
		\coordinate (A)  at (-1.5,0);
		\coordinate (a) at (-2.5,0); 
		\coordinate  (b) at (0,2.5);
		\coordinate (b')  at (0,-2.5);
		\coordinate (B) at  ($(A)+(0,1)$);
		\coordinate (B')  at ($(A')+(0,-2)$);
		 \draw[-stealth](A)--(B);
		 \draw[-stealth](A')--(B');
		\draw[stealth-stealth](b)--(b'); 
		\draw  (a)--(a')
			(B)--(B'); 
		 \draw[dashed](-1.5,0)--(-1.5,-1)
			(1,0)--(1,1.5) 
			(-1,0)--(-1,1.5); 
		 \draw[stealth-stealth,dashed](-1.5,-1)--(0,-1) node[midway,above]{$d$};
		\draw[stealth-stealth,dashed](3,-1)--(0,-1)  node[midway,above]{$d'$};
		\draw[stealth-stealth,dashed](-1,1.5)--(0,1.5) node[midway,above]{30}; 
		\draw[stealth-stealth,dashed](1,1.5)--(0,1.5)  node[midway,above]{30};
		\foreach\p/\g in {A/-45,F'/-45,F/-45,A'/-45,B/90,O/-135,B'/-90}\draw[fill=black] (\p)  circle (.1pt) node[shift={(\g:.3)},scale=1]{$\p$};
	\end{tikzpicture}
		} 
	\loigiai{Ta có $\lim\limits_{d\to {f^+}}df={f^2}$ và $\lim\limits_{d\to  {f^+}}\dfrac1{d-f}=+\infty $ nên $\lim\limits_{d\to {f^+}}\dfrac{df}{d-f}=\lim\limits_{d\to {f^+}}\left(df\cdot\dfrac1{d-f}\right)=+\infty $.}
\end{ex}

% nguồn: Ok_Toan11_TracNghiem_Lan2
\begin{ex}%[1D3V2-8]%[ma:0e10240d]
	Giả sử khoảng cách từ đỉnh của vách đá đến mặt đất là $96$ ft. Một hòn đá rơi từ đỉnh của một vách đá xuống đất, sau khoảng thời gian $t$ giây, khoảng cách  của nó so với đỉnh của vách đá là $ s(t)=16t^2$. Tại thời điểm hòn đá chạm xuống đất vận tốc của hòn đá xấp xỉ bằng
	 \choice
		 {\True $32\sqrt6$}
		{$16\sqrt6$} 
		 {0}
		 {32}
	\loigiai{Ta  có $\lim\limits_{v\to 15}f(v)=\lim\limits_{v\to 15}\dfrac{430v}{0{,}5v^2+16{,}2v+354}=\dfrac{430\cdot15}{0{,}{{5\cdot15}^2}+16{,}2\cdot15+354}\approx 9$.\\
		Hòn  đá sẽ chạm đất khi và chỉ khi $s(t)=96=16t^2\Leftrightarrow {t^2}=6\Leftrightarrow
		\hoac{&t=\sqrt6  \text{ (tm)}\\&t=-\sqrt6 \text{ (l).}}$\\
		Vậy hòn đá sẽ chạm đất tại thời điểm $  t=\sqrt6\approx 2{,}4495$ (s).\\
		Tại thời điểm hòn đá chạm xuống đất vận tốc của hòn đá là 
	$$\lim\limits_{t\to\sqrt6}\dfrac{s(t)-s\left(\sqrt6\right)}{t-\sqrt6}=\lim\limits_{t\to\sqrt6}\dfrac{16t^2-16\cdot {{\sqrt6}^2}}{t-\sqrt6}=\lim\limits_{t\to\sqrt6}16\left(t+\sqrt6\right)=32\cdot\sqrt6\approx 78{,}4.$$} 
\end{ex}

% nguồn: Ok_Toan11_TracNghiem_Lan2
\begin{ex}%[1D3V2-8]%[ma:0be282a4]
	Số lượng xe ô tô vào một đường hầm trong $1$ giây được cho bởi công thức $ f(v)=\dfrac{430v}{0{,}5v^2+16{,}2v+354}$, trong đó $v$ (m/s) là vận tốc trung bình của các xe khi đi vào đường hầm. Khi các xe vào đường hầm với vận tốc xấp xỉ $15$ (m/s) thì số xe vào  đường hầm mỗi giây là
	 \choice
		{\True $9$} 
		 {8}
		{7} 
		{6} 
	\loigiai{Ta có 
	$$\lim\limits_{v\to  15}f(v)=\lim\limits_{v\to 15}\dfrac{430v}{0{,}5v^2+16{,}2v+354}=\dfrac{430\cdot15}{0{,}5\cdot15^2+16{,}2\cdot15+354}\approx 9.$$}
\end{ex}


%%%%%%%%%%%%%%%% TRẢ LỜI NGẮN %%%%%%%%%%%%%%%%

% nguồn: Ok_TraLoiNgan_Lop11_Lan1
\begin{ex}%[1D3V2-8]%[ma:e653dd3d]
	Một cái hồ chứa $600$  lít nước ngọt. Người ta bơm nước biển có nồng độ muối $30$ gam/lít vào hồ với tốc độ $15$ lít/phút. Nồng độ muối trong hồ dần về bao nhiêu gam/lít khi $t$ dần về dương vô cùng?
	\shortans{30} 
	\loigiai{ 
		Sau $t$ phút bơm nước vào hồ thì lượng nước là $600+15t$ (lít) và lượng muối có được là $30 \cdot 15t$ (gam).\\ 
		Nồng độ muối của nước là $C(t)=\dfrac{30 \cdot 15t}{600+15t}=\dfrac{30t}{40+t}$ (gam/lít).\\ 
		 Khi $t$ dần về dương vô cùng, ta có\\
		$\lim\limits_{t\to+\infty}C(t)=\lim\limits_{t\to+\infty}\dfrac{30t}{40+t}= \lim\limits_{t\to+\infty}\dfrac{30t}{t\left(\dfrac{40}{t}+1\right)}= \lim\limits_{t\to+\infty}\dfrac{30}{\dfrac{40}{t}+1}=30$ (gam/lít). 
	} 
\end{ex}

% nguồn: Ok_TraLoiNgan_Lop11_Lan1
\begin{ex}%[1D3V2-8]%[ma:cd2ca1a8]
	Chi phí (đơn vị: nghìn đồng) để sản xuất $x$ sản phẩm của một công ty được xác định bởi hàm số $C(x)=50000+ 120x$. Biết rằng số lượng sản phẩm sản xuất càng nhiều thì chi phí sản xuất sẽ càng giảm. Khi đó chi phí sản xuất một sản phẩm  thấp nhất là bao nhiêu nghìn đồng?
	 \shortans{120}
	\loigiai{ 
		$-$ Chi phí trung bình sản xuất một sản phẩm là $\overline{C}(x)= \dfrac{50000+ 120x  }{x}$, $(x>0)$.\\
		Vì số lượng sản phẩm sản xuất càng nhiều thì chi phí sản xuất sẽ càng giảm nên chi phí sản  xuất một sản phẩm thấp nhất là\\
		$\lim\limits_{x\to  +\infty} \overline{C}(x)= \lim\limits_{x\to +\infty} \dfrac{50000+ 120x }{x}=\lim\limits_{x\to +\infty}\left(\dfrac{50000}{x}+120\right)= 120$.\\
		Vậy chi phí sản  xuất một sản phẩm thấp nhất là $120$ nghìn đồng.
	} 
\end{ex}

% nguồn: Ok_TraLoiNgan_Lop11_Lan2
\begin{ex}%[1D3V2-8]%[ma:ae0dfee2]
	Một cái hồ chứa $600$ lít nước ngọt. Người ta bơm nước biển có nồng độ muối $30$ gam/lít vào  hồ với tốc độ $15$ lít/phút. Nồng độ muối trong hồ khi $t$ dần về dương vô cùng (đơn vị: gam/lít) là bao nhiêu?
	\shortans{30} 
	\loigiai{Sau $t$ phút bơm nước vào hồ thì lượng nước là $600+15t$ (lít) và lượng muối có được là $30\cdot15t$  gam.\\
		Nồng độ muối của nước là  $C(t)=\dfrac{30\cdot15t}{600+15t}=\dfrac{30t}{40+t}$ (gam/lít).\\
		Khi $t$ dần về dương vô cùng,  ta có
		\begin{align*}
			\lim\limits_{t\to+\infty }C(t)&=\lim\limits_{t\to+\infty }\dfrac{30t}{40+t}=\lim\limits_{t\to+\infty }\dfrac{30t}{t\left(\dfrac{40}t+1\right)}=\lim\limits_{t\to+\infty  }\dfrac{30}{\dfrac{40}t+1}=30 \text{ (gam/lít).}
	\end{align*}}
\end{ex}

% nguồn: Ok_TraLoiNgan_Lop11_Lan1
\begin{ex}%[1D3C2-8]%[ma:06e9bf10]
	 \immini
	{Trong hệ trục toạ độ $Oxy$, lấy điểm $A$ thuộc tia $Ox$ và điểm $B(0;2)$ thuộc tia $Oy$. Giả sử hoành độ điểm $A$ là $a > 0$. Khi điểm $A$ dịch chuyển ra vô cực theo chiều dương trục  $Ox$ thì độ dài $AH$ dần về bao nhiêu?}
	{ 
	\begin{tikzpicture}[scale=1,line join=round,line cap=round,font=\footnotesize,>=stealth]
			\path (0,0) coordinate (O); 
			\draw[-stealth] (0,0)--(4,0) node[above]{$x$}; 
			\draw[-stealth] (0,0)--(0,3) node[left]{$y$}; 
			\node[left] at (0,2) {$2$}; 
			\node[below] at  (3,0) {$a$};
			 \node[right] at (0,2) {$B$};
			 \node[above] at (3,0) {$A$};
			\node[above] at (1,1.5) {$H$}; 
			\draw (0,2)--(3,0) 
				 (0,0)--(0.9,1.4);
		\end{tikzpicture}
	 }
	 \shortans{2}
	 \loigiai{
		 Tam giác $OAB$ vuông tại $O$ có đường cao $OH$ nên:\\
		 $\dfrac{1}{O{H^2}}=\dfrac{1}{O{A^2}}+\dfrac{1}{O{B^2}}\Rightarrow OH= \dfrac{OA\cdot OB}{\sqrt{O{A^2}+O{B^2}}}= \dfrac{2a}{\sqrt{4+a^2}}$.\\
		Đặt  $h(a)=OH=\dfrac{2a}{\sqrt{4+a^2}}$.\\
		Khi  điểm $A$ dịch chuyển ra vô cực theo tia $Ox$ thì $a\to+\infty $.\\
		Ta có $\lim\limits_{a\to+\infty}h(a)= \lim\limits_{a\to+\infty}\dfrac{2a}{\sqrt{4+a^2}}=  \lim\limits_{a\to+\infty}\dfrac{2a}{a\sqrt{\dfrac{4}{a^2}+1}}= \lim\limits_{a\to+\infty}\dfrac{2}{\sqrt{\dfrac{4}{a^2}+1}}= \dfrac{2}{\sqrt 1}=2$.\\
		 Vậy khi điểm $A$ dần về vô cực thì độ dài $OH$ dần về $2$.
	 }
\end{ex}


%%%%%%%%%%%%%%%% TỰ LUẬN %%%%%%%%%%%%%%%%

% nguồn: Ok_TuanLuan_Toan11_Lan2
\begin{ex}%[1D3V2-8]%[ma:37f74c91]

	Chi phí  để sản xuất $x$ chai nước ngọt của công ty nước giải khát A được xác định bởi hàm số $F(x) = 50000 + 15x$ (đơn vị: nghìn đồng).
	\begin{enumerate}
		\item Tính chi phí trung bình $\overline{F}(x)$ để công ty  sản xuất một sản phẩm.
		\item Tính $\displaystyle\lim_{x \to +\infty}  \overline{F}(x)$ và cho biết ý nghĩa của kết quả.
	\end{enumerate}
	 \loigiai{
		\begin{enumerate}
			\item Chi  phí trung bình để sản xuất $x$ sản phẩm là
			$\overline{F}(x) = \dfrac{50000 +  15x}{x}$ (nghìn đồng).
			\item  Ta có
			\begin{eqnarray*}
				&&\displaystyle\lim_{x  \to +\infty} \overline{F}(x) \\
				&=& \displaystyle\lim_{x  \to +\infty} \dfrac{50000 + 15x}{x}\\
				&=& \displaystyle\lim_{x \to +\infty} \left( \dfrac{50000}{x} + 15 \right)\\ 
				&=&  15.
			\end{eqnarray*}
			\textbf{Ý nghĩa:} Khi số sản phẩm công ty sản xuất ra càng nhiều thì chi phí trung bình để sản xuất ra một sản  phẩm sẽ giảm dần, số sản phẩm $x$ đủ lớn thì chi phí trung bình xấp xỉ $15$ nghìn đồng mỗi sản phẩm và không thể thấp hơn.
		\end{enumerate}
	 }
\end{ex}

% nguồn: Ok_TuanLuan_Toan11_Lan2
\begin{ex}%[1D3V2-8]%[ma:51b13b63]

	Một  hồ nuôi tôm chứa $600$ m$^3$ nước mặn với nồng độ muối $1$ kg/m$^3$. Chủ hồ nuôi tôm dự định chuyển đổi giống mới nên bơm nước vào hồ với vận tốc $3$ m$^3$/phút để làm ngọt hóa nước trong hồ.
	\begin{enumerate}
		\item Viết biểu thức $C(x)$ biểu thị nồng độ muối trong hồ sau $x$ phút kể từ khi  bắt đầu bơm.
		\item Tính $\displaystyle\lim_{x \to +\infty} C(x)$ và  giải thích ý nghĩa của kết quả này.
	\end{enumerate}
	 \loigiai{
		\begin{enumerate}
			\item Khối lượng muối có  trong hồ nuôi tôm là $1\cdot 600= 600$ (kg).\\
			Sau $x$ phút, lượng nước  trong hồ là $600 + 3x$ (m$^3$).\\
			Nồng độ muối tại thời  điểm $x$ phút kể từ khi bơm thêm nước ngọt vào là
			$$ C(x) =  \dfrac{600}{600 + 3x}.$$
			\item Ta  có
			\begin{eqnarray*}
				&&\displaystyle\lim_{x \to +\infty} C(x)\\ 
				&=& \displaystyle\lim_{x \to +\infty} \dfrac{600}{600  + 3x}\\
				 &=& 0.
			\end{eqnarray*}
			\textbf{Ý nghĩa:} Điều này có nghĩa là khi $x$ càng lớn thì nồng độ muối trong hồ sẽ dần về  $0$. Tức là đến một thời điểm nào đó muối trong hồ không còn đáng kể và nước trong hồ coi như nước ngọt.
		\end{enumerate}
	 }
\end{ex}

% nguồn: Ok_TuanLuan_Toan11_Lan2
\begin{ex}%[1D3V2-8]%[ma:14e2e768]
	Cho hàm số $f(x)$ xác  định trên $\mathbb{R}$ thỏa mãn $\lim\limits_{x\to 2} \dfrac{f(x)-16}{x-2}=12$. Giới hạn $\lim\limits_{x\to 2} \dfrac{\sqrt[3]{5f(x)-16}-4}{x^2+2x-8}$ bằng
	\loigiai{ 
		Do  $\lim\limits_{x\to 2}\dfrac{f(x)-16}{x-2}=12$ nên ta có $f(2)-16=0$ hay $f(2)=16$.
		 \allowdisplaybreaks
		\begin{eqnarray*}
			\lim\limits_{x\to  2} \dfrac{\sqrt[3]{5f(x)-16}-4}{x^2+2x-8}&=&\lim\limits_{x\to 2} \dfrac{5(f(x)-16)}{(x-2)(x+4)(\sqrt[3]{(5f(x)-16)^2}+4\sqrt[3]{5f(x)-16}+16)}\\
			&=&\lim\limits_{x\to  2} \dfrac{f(x)-16}{x-2}\cdot \dfrac{5}{(x+4)(\sqrt[3]{(5f(x)-16)^2}+4\sqrt[3]{5f(x)-16}+16)}\\
			&=&12\cdot  \dfrac{5}{6\cdot 48}=\dfrac{5}{24}.
		\end{eqnarray*}
	} 
\end{ex}

% nguồn: Ok_TuanLuan_Toan11_Lan2
\begin{ex}%[1D3V2-8]%[ma:d1fd0dba]
	Một công ty sản xuất máy tính cầm tay đã xác định được rằng,  tính trung bình một nhân viên có thể lắp ráp được $C(t)=\dfrac{30t}{5+t}(t\ge 0)$ chiếc máy tính mỗi ngày sau khi được đào tạo $t$ ngày. Tính $\lim\limits_{t\to +\infty}C(t)$ và cho biết ý nghĩa của kết quả này.
	 \loigiai{
		 Ta có $\lim\limits_{t\to +\infty}C(t)
		=\lim\limits_{t\to +\infty}\dfrac{30t}{5+t} 
		 =\lim\limits_{t\to +\infty}\left(\dfrac{30}{\dfrac{5}{t}+1}\right)=30$.\\
		 \textbf{\textit{Ý nghĩa:}} Khi thời gian đào tạo một nhân viên tăng lên thì số máy tính trung bình mỗi nhân viên lắp ráp được trong một ngày cũng tăng lên, như không quá $30$ chiếc mỗi ngày.
	 }
\end{ex}

% nguồn: Ok_TuanLuan_Toan11_Lan2
\begin{ex}%[1D3V2-8]%[ma:5ad0b7ff]

	Trong thuyết tương đối của Einstein, khối lượng của vật chuyển động với vận tốc $v$ cho bởi công thức $m = \dfrac{m_0}{\sqrt{1 - \dfrac{v^2}{c^2}}}$, trong đó  $m_0$ là khối lượng của vật khi nó đứng yên, $c$ là vận tốc ánh sáng. Chuyện gì xảy ra với khối lượng của vật khi vận tốc của vật gần với vận tốc ánh sáng?
	\loigiai{ 
		Từ công  thức khối lượng $m = \dfrac{m_0}{\sqrt{1 - \dfrac{v^2}{c^2}}}$, ta thấy $m$ là một hàm số của $v$, với tập xác định là nửa khoảng $\left[0;c\right)$.\\
		 Rõ ràng khi $v$ tiến gần tới vận tốc ánh sáng, tức là $v \to c^-$, ta có $\sqrt{1-\dfrac{v^2}{c^2}} \to 0$.\\
		Do đó $\displaystyle\lim_{v \to c^-} m\left(v\right)=+\infty$, nghĩa là khối lượng $m$ của vật trở nên vô cùng lớn khi vận tốc của vật gần với  vận tốc ánh sáng.
	 }
\end{ex}

% nguồn: Ok_TuanLuan_Toan11_Lan2
\begin{ex}%[1D3V2-8]%[ma:f47552d1]
	Cho đa thức $f(x)$ thỏa  mãn $\lim\limits_{x\to 1}\dfrac{f(x)-2}{x-1}=12$. Tính $\lim\limits_{x\to 1} \dfrac{f(x)-2}{\left(x^2-1\right)\left[f(x)+1\right]}$
	\loigiai{ 
		\begin{itemize}
			\item Cách 1.\\ 
			Theo giả thiết bài toán ta có  $\lim\limits_{x\to 1} \left[f(x)-2\right]=0\Rightarrow \lim\limits_{x\to 1}f(x)=2$.\\
			Khi  đó
			\begin{eqnarray*}
				 \lim\limits_{x\to 1} \dfrac{f(x)-2}{\left(x^2-1\right)\left[f(x)+1\right]}&=&\lim\limits_{x\to 1} \dfrac{f(x)-2}{(x-1)(x+1)\left[f(x)+1\right]}\\ &=&12\lim\limits_{x\to 1} \dfrac{1}{(x+1)\left[f(x)+1\right]}=12\cdot \dfrac{1}{2\cdot (2+1)}=\dfrac{12}{6}=2.
			\end{eqnarray*}
			\item Cách  2.\\
			Giả sử $f(x)=12(x-1)+2$ thỏa mãn  giả thiết bài toán $\lim\limits_{x\to 1}\dfrac{f(x)-2}{x-1}=12$.\\
			Khi đó  $\lim\limits_{x\to 1} \dfrac{f(x)-2}{\left(x^2-1\right)\left[f(x)+1\right]}=\lim\limits_{x\to 1}\dfrac{12(x-1)}{(x-1)(x+1)(12(x-1)+3)}=\dfrac{12}{6}=2$.
		\end{itemize}
	 }
\end{ex}

% nguồn: Ok_TuanLuan_Toan11_Lan2
\begin{ex}%[1D3V2-8]%[ma:147db317]
	Cho hàm số $f(x)$ xác định trên $\mathbb{R}$ thỏa mãn $\lim\limits_{x\to 2} \dfrac{f(x)-16}{x-2}=12$. Tính giới hạn\\ $\lim\limits_{x\to 2}  \dfrac{\sqrt{2f(x)-16}-4}{x^2+x-6}$.
	 \loigiai{
		Vì $\lim\limits_{x\to 2}\dfrac{f(x)-16}{x-2}=12$ nên $\lim\limits_{x\to 2}\left[f(x)-16\right]=0$.\\ 
		Do đó nếu giới hạn này khác $0$ thì giới hạn $\lim\limits_{x\to 2} \dfrac{f(x)-16}{x-2}$ sẽ bằng vô cùng.\\ Ta  suy ra được $\lim\limits_{x\to 2}f(x)=16$.\\
		 Biến đổi
		\allowdisplaybreaks 
		\begin{eqnarray*}
			\lim\limits_{x\to 2} \dfrac{\sqrt{2f(x)-16}-4}{x^2+x-6}&=&\lim\limits_{x\to 2} \dfrac{2f(x)-32}{(x-2)(x+3)\left(\sqrt{2f(x)-16}+4\right)}\\ &=&\lim\limits_{x\to 2}\left[\dfrac{f(x)-16}{\left(x-2\right)}\cdot \dfrac{2}{(x+3)\left(\sqrt{2f(x)-16}+4\right)}\right] 
		\end{eqnarray*}
		Do  $\lim\limits_{x\to 2}f(x)=16$ nên $\lim\limits_{x\to 2} \left[\dfrac{2}{(x+3)\left(\sqrt{2f(x)-16}+4\right)} \right]=\dfrac{1}{20}$.\\
		Vậy $\lim\limits_{x\to 2} \dfrac{\sqrt{2f(x)-16}-4}{x^2+x-6}=\lim\limits_{x\to 2} \left[\dfrac{f(x)-16}{(x-2)}\cdot  \dfrac{2}{(x+3)\left(\sqrt{2f(x)-16}+4\right)}\right]=12\cdot \dfrac{1}{20}=\dfrac{3}{5}$.
	} 
\end{ex}

% nguồn: Ok_TuanLuan_Toan11_Lan2
\begin{ex}%[1D3V2-8]%[ma:c392497a]
	Số lượng xe ô tô vào một đường hầm trong $1$ giây được cho bởi công thức sau đây: $f(v)=\dfrac{290{,}4v}{0{,}36v^2+13{,}2v+264}$, trong đó $v$ (m/s) là vận  tốc trung bình của các xe khi đi vào đường hầm $\lim\limits_{v\to 15}f(v)$ và cho biết ý nghĩa của kết quả (làm tròn kết quả đến hàng đơn vị).
	 \loigiai{
		Ta có  $\lim\limits_{v\to 15}f(v)
		=\lim\limits_{v\to 15}\dfrac{290{,}4v}{0{,}36v^2+13{,}2v+264} 
		=\dfrac{290{,}4\cdot 15}{0{,}36\cdot  15^2+13{,}2\cdot 15+264}\approx 8$.\\
		Ý nghĩa: Khi các xe vào đường hầm với vận tốc  xấp xỉ $20$ (m/s) thì trung bình mỗi giây có khoảng $8$ xe vào đường hầm.
	 }
\end{ex}

% nguồn: Ok_TuanLuan_Toan11_Lan2
\begin{ex}%[1D3V2-8]%[ma:25b84aaa]
	Một bể chứa $5\,000$ lít nước tinh khiết. Nước muối có chứa $30$ gam muối trên mỗi  lít nước được bơm vào bể với tốc độ $25$ lít/phút.
	\begin{enumerate}
		\item Chứng minh rằng nồng độ muối của nước  trong bể sau $t$ phút (tính bằng khối lượng muối chia thể tích nước trong bể, đơn vị g/l) là $C(t)=\dfrac{30t}{200+t}$.
		\item Tính  $\lim\limits_{t\to +\infty}C(t)$ và cho biết ý nghĩa của kết quả đó.
	\end{enumerate}
	 \loigiai{
		\begin{enumerate}
			\item Sau $t$ phút thì lượng muối  trong bể là $30\cdot 25\cdot t=750t$ (g) và thể tích nước trong bể là $5\,000+25t$ (l).\\
			Vậy nồng độ muối của nước trong bể sau $t$ phút là $$C(t)=\dfrac{750t}{5000+25t}=\dfrac{30t}{200+t} \, \mathrm{(g/l)}.$$ 
			\item Ta  có
			 $\lim\limits_{t\to +\infty}C(t)
			=\lim\limits_{t\to  +\infty}\dfrac{30t}{200+t}
			 =\lim\limits_{t\to +\infty}\dfrac{30}{\dfrac{200}{t}+1}
			=\dfrac{30}{1}=30$.\\ 
			\textbf{\textit{Ý nghĩa:}} Theo kết quả đó, ta thấy khi lượng nước trong bể tăng theo thời gian  đến vô hạn thì nồng độ muối của nước sẽ tăng dần đến giá trị $30$ (g/l), tức là xấp xỉ nồng độ muối của loại nước muối cho thêm vào bể.
		\end{enumerate}
	 }
\end{ex}

% nguồn: Ok_TuanLuan_Toan11_Lan2
\begin{ex}%[1D3V2-8]%[ma:ec6f59a7]
	Sau khi phát hiện một bệnh dịch, các chuyên gia y tế ước tính số người nhiễm bệnh kể từ ngày xuất hiện bệnh nhân đầu tiên biến đổi theo một hàm số thời gian (tính theo ngày) là  $g(t)=45t^2-t^3$ (người). Tốc độ trung bình gia tặng người bệnh giữa hai thời điểm $t_1$, $t_2$ là $v_{\text{tb}}=\dfrac{g(t_2)-g(t_1)}{t_2-t_1}$. Tính $\lim\limits_{t\to 10}\dfrac{g(t)-g(10)}{t-10}$ và cho biết ý nghĩa của kết quả tìm được.
	 \loigiai{
		Ta  có $g(10)=45\cdot 10^2-10^3$.\\
		Khi đó 
		\begin{eqnarray*}
			 \lim\limits_{t\to 10}\dfrac{g(t)-g(10)}{t-10}
			&=&\lim\limits_{t\to 10}\dfrac{45t^2-t^3-45\cdot 10^2-10^3}{t-10} \\ 
			&=&\lim\limits_{t\to 10}\dfrac{45(t^2-10^2)-(t^3-10^3)}{t-10}  \\
			&=&\lim\limits_{t\to 10}\dfrac{45(t-10)(t+10)-(t-10)(t^2+10t+100)}{t-10}\\ 
			&=&\lim\limits_{t\to 10}\left[45(t+10)-(t^2+10t+100)\right]\\  &=&\lim\limits_{t\to 10}(-t^2+35t+350)=600.
		\end{eqnarray*}
		\textbf{\textit{Ý nghĩa:}}  Từ kết quả trên, ta thấy tốc độ tăng người bệnh ngay tại thời điểm $t=10$ ngày là $600$ người/ngày.
	 }
\end{ex}

% nguồn: Ok_TuanLuan_Toan11_Lan2
\begin{ex}%[1D3V2-8]%[ma:22e06722]
	Một hồ nước chứa $500$ m$^3$ nước mặn với nồng độ muối $6$ kg/m$^3$. Người ta dự định ngọt hóa nước trong hồ bằng cách bơm nước vào  hồ với vận tốc $3$ (m$^3$/ phút).
	\begin{enumerate}
		\item Viết biểu  thức $C(t)$ biểu thị nồng độ muối trong hồ sau $t$ phút kể từ khi bắt đầu bơm.
		\item Tìm giới hạn $\lim\limits_{t\to +\infty}C(t)$ và giải thích  ý nghĩa của kết quả này.
	\end{enumerate}
	 \loigiai{
		\begin{enumerate}
			\item Khối lượng muối có trong hồ  là $500\cdot 6=3\,000(kg)$.\\
			Sau $t$ phút,  lượng nước trong hồ là $500+3t$ (m$^3$).\\
			Nồng độ muối tại thời điểm $t$ phút kể từ khi bơm là  $C(t)=\dfrac{3000}{500+3t}$ (kg/m$^3$).
			\item Ta có $\lim\limits_{t\to +\infty}C(t) 
			=\lim\limits_{t\to  +\infty}\dfrac{3\,000}{500+3t}
			 =\lim\limits_{t\to +\infty}\dfrac{\dfrac{3\,000}{t}}{\dfrac{500}{t}+3}
			=\dfrac{0}{3}=0$. 
		\end{enumerate}
		Điều này có nghĩa là khi $t$ càng lớn thì nồng độ muối trong hồ sẽ dần về $0$. Tức là, đến một thời điểm nào đó nồng độ muối trong hồ không còn đáng kể và nước  trong hồ xem như nước ngọt.
	 }
\end{ex}

% nguồn: Ok_TuanLuan_Toan11_Lan2
\begin{ex}%[1D3V2-8]%[ma:d7adf609]
	Cho  hàm số $y=f(x)$ có giới hạn hữu hạn tại $x=4$ và thỏa mãn $\lim\limits_{x\to 4}\dfrac{f(x)+3}{x-4}=1$.\\ Tính $L=\lim\limits_{x\to 4} \dfrac{\sqrt[3]{f(x)+4}-1}{x-4}$.
	 \loigiai{
		Nếu $\lim\limits_{x\to 4}f(x)\ne-3$ thì $\lim\limits_{x\to 4} \left[f(x)+3\right]\ne 0$ và  $\lim\limits_{x\to 4}(x-4)=0$.\\
		Khi đó  $\lim\limits_{x\to 4}\dfrac{f(x)+3}{x-4}$ không thể hữu hạn. Do đó, $\lim\limits_{x\to 4}f(x)=-3$.\\
		 Ta có
		 \allowdisplaybreaks
		\begin{eqnarray*}
			L&=& \lim\limits_{x\to 4}  \dfrac{\sqrt[3]{f(x)+4}-1}{x-4}\\
			&=&  \lim\limits_{x\to 4} \dfrac{f(x)+3}{\left(x-4\right)\left[\sqrt[3]{\left(f(x)+4\right)^2}+\sqrt[3]{f(x)+4}+1\right]}\\
			&=& \lim\limits_{x\to 4}\left[\dfrac{f(x)+3}{x-4}\cdot \dfrac{1}{\sqrt[3]{\left(f(x)+4\right)^2}+\sqrt[3]{f(x)+4}+1}  \right]\\
			 &=& 1\cdot \dfrac{1}{3}\\
			&=& \dfrac{1}{3}. 
		\end{eqnarray*}
	} 
\end{ex}

% nguồn: Ok_TuanLuan_Toan11_Lan2
\begin{ex}%[1D3V2-8]%[ma:57f78b23]
	Chi phí (đơn vị: nghìn đồng) để sản  xuất $x$ sản phẩm của một công ty được xác định bởi hàm số $F(x)=60\,000+250x$.
	\begin{enumerate}
		\item Tính chi phí trung bình $\overline{F}(x)$ để  công ty sản xuất một sản phẩm.
		\item Tính $\lim\limits_{x\to +\infty}\overline{F}(x)$ và cho biết ý nghĩa của kết  quả.
	\end{enumerate}
	 \loigiai{
		\begin{enumerate}
			\item Chi phí trung bình để sản  xuất $x$ sản phẩm là $\overline{F}(x)=\dfrac{60\,000+250x}{x}$ (nghìn đồng).
			\item Ta  có $\lim\limits_{x\to +\infty}\overline{F}(x)
			=\lim\limits_{x\to  +\infty}\dfrac{60\,000+250x}{x}
			=\lim\limits_{x\to +\infty}\left(\dfrac{60\,000}{x}+250\right)=250$ (nghìn  đồng).
		\end{enumerate}
		\textbf{\textit{Ý nghĩa:}} Khi số sản phẩm công ty sản xuất ra càng nhiều thì chi phí trung bình để sản xuất ra một sản phẩm sẽ giảm dần, số sản phẩm $x$ đủ lớn thì  chi phí trung bình xấp xỉ $250$ (nghìn đồng) mỗi sản phẩm và không thể thấp hơn.
	 }
\end{ex}

% nguồn: Ok_TuanLuan_Toan11_Lan2
\begin{ex}%[1D3V2-8]%[ma:8c919994]

	Số lượng xe ô tô  vào một đường hầm Thủ Thiêm trong một giây được cho bởi công thức $f\left(v\right) = \dfrac{400v}{v^2+3v+6}$, trong đó $v$ (m/s) là vận tốc trung bình của các xe khi đi vào đường hầm. Tính $\displaystyle\lim_{v \to 20} f\left(v\right)$ và cho biết ý nghĩa của kết quả (làm tròn kết quả đến hàng đơn vị).
	 \loigiai{
		 Ta có
		\begin{eqnarray*}
			&&\displaystyle\lim_{v \to 20} f\left(v\right)\\ 
			 &=& \displaystyle\lim_{v \to 20} \dfrac{400v}{v^2+3v+6}\\
			&=& \frac{8000}{466}  \approx 17.
		\end{eqnarray*}
		\textbf{Ý nghĩa:} Khi các xe vào đường hầm với vận tốc xấp xỉ $20$ m/s thì trung bình mỗi giây có khoảng $17$ xe vào  đường hầm.
	} 
\end{ex}

% nguồn: Ok_TuanLuan_Toan11_Lan2
\begin{ex}%[1D3V2-8]%[ma:ce09987a]

	Một công ty sản  xuất giày da đã xác định được rằng, tính trung bình một công nhân có thể làm được $f(x) = \dfrac{16x}{15 + 2x}$ (đôi giày) mỗi ngày sau khi được đào tạo $x$ ngày. Tính $\displaystyle\lim_{x \to +\infty} f(x)$ và giải thích ý nghĩa của kết quả này.
	\loigiai{ 
		Ta có 
		\begin{eqnarray*}
			 &&\displaystyle\lim_{x \to +\infty} f(x)\\
			 &=& \displaystyle\lim_{x \to +\infty} \dfrac{16x}{15 + 2x}\\
			&=&  8.
		\end{eqnarray*}
		\textbf{Ý nghĩa:} Khi thời gian đào tạo tăng lên  thì số đôi giày mỗi công nhân sản xuất được trong một ngày cũng được củng cố tăng lên nhưng không quá $8$ đôi giày/ngày.
	} 
\end{ex}

% nguồn: Ok_TuanLuan_Toan11_Lan2
\begin{ex}%[1D3V2-8]%[ma:08e10641]

	Một  bar ngoài trời ở khu du lịch A đã xác định rằng, tính trung bình một ngày số lượng khách du lịch đến bar được xác định bởi công thức $f\left(x\right) = \dfrac{6000x}{50 + 12x}$ khách mỗi ngày sau khi số nhân viên của bar được tăng lên $x$ nhân viên. Tính $\displaystyle\lim_{x \to +\infty} f\left(x\right)$ và giải thích ý nghĩa của kết quả này.
	\loigiai{ 
		 Ta có
		\begin{eqnarray*}
			 &&\displaystyle\lim_{x \to +\infty} f\left(x\right)\\
			&=& \displaystyle\lim_{x \to +\infty} \dfrac{6000x}{50  + 12x}\\
			&=& \displaystyle\lim_{x \to +\infty} \dfrac{6000}{50/x  + 12}\\
			&=& 500. 
		\end{eqnarray*}
		\textbf{Ý nghĩa:} Khi số nhân viên trong bar tăng lên thì số khách đến bar trong một một  ngày cũng được cũng tăng lên nhưng không quá $500$ khách/ ngày.
	 }
\end{ex}

% nguồn: Ok_TuanLuan_Toan11_Lan2
\begin{ex}%[1D3V2-8]%[ma:a2a82949]

	Chi phí để sản xuất $x$ kg chả mực của cửa hàng  bán chả mực A được xác định bởi hàm số $F\left(x\right)=60000+320x$ (đơn vị: nghìn đồng).
	\begin{enumerate}
		\item Tính chi phí trung bình $\overline{F}\left(x\right)$ để công ty sản xuất  $x$ kg chả mực.
		\item Tính $\displaystyle\lim_{x \to +\infty} \overline{F}\left(x\right)$ và cho biết ý nghĩa  của kết quả.
	\end{enumerate}
	\loigiai{ 
		\begin{enumerate}
			 \item Chi phí trung bình để sản xuất $x$ sản phẩm là $\overline{F}\left(x\right)=\dfrac{60000+320x}{x}$ (nghìn đồng).
			 \item Ta có
			\begin{eqnarray*}
				 &&\displaystyle\lim_{x \to +\infty} \overline{F}\left(x\right)\\
				&=& \displaystyle\lim_{x \to +\infty}  \dfrac{60000+320x}{x}\\
				&=& \displaystyle\lim_{x\to +\infty}  \left(\dfrac{60000}{x}+320\right)\\
				&=&320. 
			\end{eqnarray*}
			\textbf{Ý nghĩa:} Khi khối lượng chả mực cửa  hàng sản xuất ra càng nhiều thì chi phí trung bình để sản xuất ra một kg sản phẩm sẽ giảm dần, số sản phẩm $x$ đủ lớn thì chi phí trung bình xấp xỉ $320$ nghìn đồng mỗi kg chả mực và không thể thấp hơn.
		\end{enumerate}
	} 
\end{ex}

% nguồn: Ok_TuanLuan_Toan11_Lan2
\begin{ex}%[1D3V2-8]%[ma:be2efa1b]
	Một thấu kính hội tụ có tiêu cự là $f$. Gọi $d$ và $d'$ lần lượt là khoảng cách  từ một vật thật $AB$ và từ ảnh $A'B'$ của nó tới quang tâm $O$ của thấu kính như hình vẽ. Công thức thấu kính là $\dfrac{1}{d}+\dfrac{1}{d'}=\dfrac{1}{f}$.
	\begin{enumerate}
		\item Tìm biểu  thức xác định hàm số $d'=\varphi (d)$.
		\item Tìm $\lim\limits_{d\to f^+} \varphi (d)$, $\lim\limits_{d\to  f^-}\varphi (d)$ và $\lim\limits_{d\to f}\varphi (d)$. Giải thích ý nghĩa của các kết quả tìm được.
	\end{enumerate}
	 \loigiai{
		\begin{enumerate}
			\item Theo giả  thiết: $\dfrac{1}{d'}=\dfrac{1}{f}-\dfrac{1}{d}
			\Leftrightarrow \dfrac{1}{d'}=\dfrac{d-f}{df} 
			\Leftrightarrow d'=\dfrac{df}{d-f}$. 
			\item Ta có 
			$\lim\limits_{d\to f^+}\varphi  (d)=\lim\limits_{d\to f^+}\dfrac{df}{d-f}=+\infty$; $\lim\limits_{d\to f^-}\varphi (d)=\lim\limits_{d\to f^-}\dfrac{df}{d-f}=-\infty$.\\
			Và $\lim\limits_{d\to f}\varphi (d)=\lim\limits_{d\to  f}\dfrac{df}{d-f}=\infty$.\\
			\textbf{\textit{Ý nghĩa:}} Khi khoảng cách của vật tới  thấu kính mà gần với tiêu cự thì khoảng cách ảnh của vật đến thấu kính ra xa vô tận nên lúc đó bằng mắt thường mình không nhìn thấy.
		\end{enumerate}
	} 
\end{ex}

% nguồn: Ok_TuanLuan_Toan11_Lan2
\begin{ex}%[1D3V2-8]%[ma:cf69e5a7]
	Cho hàm số $y=f(x)$ có giới hạn hữu hạn tại $x=-1$  và thỏa mãn $\lim\limits_{x\to-1}\dfrac{f(x)+3}{x^2-1}=2$. Tính $\lim\limits_{x\to-1} \dfrac{\sqrt{f(x)+7}-\sqrt[3]{2f(x)+14}}{x^3+x^2}$.
	\loigiai{ 
		Nếu $\lim\limits_{x\to-1}f(x)\ne-3$ thì $\lim\limits_{x\to-1}\left[f(x)+3\right]\ne 0$ và  $\lim\limits_{x\to-1}\left(x^2-1\right)=0$.\\
		Khi đó $\lim\limits_{x\to-1}\dfrac{f(x)+3}{x^2-1}$ không thể hữu hạn. Do đó, $\lim\limits_{x\to-1}  f(x)=-3$.\\
		 Ta có
		 \allowdisplaybreaks
		\begin{eqnarray*}
			A&=&\lim\limits_{x\to-1}  \dfrac{\sqrt{f(x)+7}-\sqrt[3]{2f(x)+14}}{x^3+x^2}\\
			&=&\lim\limits_{x\to-1} \left[\dfrac{\sqrt{f(x)+7}-2}{x^3+x^2}-\dfrac{\sqrt[3]{2f(x)+14}-2}{x^3+x^2}  \right]\\
			 &=&\lim\limits_{x\to-1} \left(\dfrac{f(x)+3}{x^2\left(x+1\right)\left[\sqrt{f(x)+7}+2\right]}-\dfrac{2f(x)+6}{x^2\left(x+1\right)\left[\sqrt[3]{\left(2f(x)+14\right)^2}+2\sqrt[3]{2f(x)+14}+4\right]} \right)\\
			&=&\lim\limits_{x\to-1} \left(\dfrac{f(x)+3}{x^2-1}\cdot\left[\dfrac{x-1}{x^2.\left(\sqrt{f(x)+7}+2\right)}-\dfrac{2\left(x-1\right)}{x^2\cdot \left(\sqrt[3]{\left(2f(x)+14\right)^2}+2\sqrt[3]{2f(x)+14}+4\right)} \right]\right)\\ 
			 &=&2\cdot \left(\dfrac{-2}{2+2}-\dfrac{-4}{4+4+4} \right)=-\dfrac{1}{3}.
		\end{eqnarray*}
	 }
\end{ex}

% nguồn: Ok_TuanLuan_Toan11_Lan2
\begin{ex}%[1D3V2-8]%[ma:8ad13fc5]
	Cho hàm số $f(x)$ xác định trên $\mathbb{R}$ thỏa mãn $\lim\limits_{x\to 2} \dfrac{f(x)-16}{x-2}=12$. Giới  hạn $\lim\limits_{x\to 2} \dfrac{\sqrt{2f(x)-16}-4}{x^2+x-6}$ bằng
	\loigiai{ 
		\begin{itemize}
			\item Vì $\lim\limits_{x\to  2}\dfrac{f(x)-16}{x-2}=12\Rightarrow f(2)=16$.
			 \item Có
			\begin{eqnarray*}
				\lim\limits_{x\to 2}  \dfrac{\sqrt{2f(x)-16}-4}{x^2+x-6}&=&\lim\limits_{x\to 2} \dfrac{2f(x)-32}{\left(x^2+x-6\right)\left(\sqrt{2f(x)-16}+4\right)}\\
				 &=&\lim\limits_{x\to 2}\dfrac{f(x)-16}{x-2}\cdot \dfrac{2}{(x+3)\left(\sqrt{2f(x)-16}+4\right)}\\
				&=&12\cdot  \dfrac{2}{(2+3)\left(\sqrt{2f(2)-16}+4\right)}\\
				&=&12\cdot \dfrac{2}{5\left(\sqrt{2\cdot  16-16}+4\right)}=\dfrac{3}{5}.
			\end{eqnarray*}
		\end{itemize}
	} 
\end{ex}
